% Discussion
%
% Interpret findings, connect back to RQs, discuss implications for
% researchers, practitioners, and tool vendors.
This section interprets the results reported in
Section~\ref{sec:results} against the four research questions, then
draws out implications for three audiences -- researchers, tool
vendors, and practitioners adopting these techniques in industrial
pipelines -- before turning to the quality engineering safety synthesis
that spans all four RQs.

\subsection{Revisiting RQ1: What Techniques Exist, and How Mature Are They?}
The 18-category taxonomy shows that ``next-generation software
testing'' is not a single technique but a loosely federated family of
architectures unified mainly by their reliance on an LLM as a
reasoning component somewhere in the pipeline. At one extreme sit
single-shot code-completion-style test generators that differ from
classical automated test generation only in using a pretrained
language model instead of a symbolic or search-based engine; at the
other extreme sit genuinely agentic systems that plan a multi-step
investigation, call external tools (compilers, debuggers, browsers,
issue trackers), observe the results, and revise their own plan with
little or no human supervision. The evidence base is considerably
stronger for the former than the latter: unit-test generation and
mutation-testing techniques (categories well represented in the
corpus) benefit from a decade of prior automated-test-generation
research to build on and compare against, whereas multi-agent and
GUI/mobile testing agents are newer, smaller in study count, and more
often evaluated on bespoke benchmarks constructed by the same team
that built the tool -- a pattern that, as discussed under RQ3 in
Section~\ref{sec:results}, correlates with inflated effectiveness claims. The practical
implication for RQ1 is that technique maturity should not be inferred
from architectural sophistication: a simple, well-evaluated single-shot
generator may currently be a safer adoption choice than an elaborate
multi-agent pipeline whose claimed gains have not been independently
reproduced.

\subsection{Revisiting RQ2: What Autonomy Levels Are Actually Achieved?}
Despite frequent use of the word ``autonomous'' in paper titles, the
corpus shows that most deployed systems operate with a human firmly in
the loop at one or more critical junctures: approving generated test
code before it enters a suite, triaging flagged flaky tests rather than
having them auto-quarantined, or reviewing agent-authored pull requests
before merge. Genuinely low-supervision operation is concentrated in
narrow, well-bounded sub-tasks (e.g., locator self-healing, mutation
generation) rather than across the full test life cycle (design,
generation, execution, oracle judgment, and maintenance). This gap
between rhetorical and demonstrated autonomy is precisely why this
review adopts an explicit autonomy-level framing (L0--L5, following
\cite{feng2025autonomylevels}) in the roadmap and governance discussion
below: without a shared vocabulary for \emph{how much} supervision a
system actually requires, ``autonomous testing'' risks becoming a
marketing label rather than a measurable engineering property.

\subsection{Revisiting RQ3: How Rigorous Is the Supporting Evidence?}
The clearest cross-cutting finding under RQ3 is that benchmark design
problems are not isolated incidents but a structural feature of this
literature. Studies in the ``Benchmarks and data contamination''
category document cases where headline accuracy numbers for neural
test-oracle generation and flaky-test classification fell substantially
once training/test leakage and class-balance artifacts were corrected,
and several comparative studies found that simple, non-agentic
baselines matched or exceeded elaborate multi-agent pipelines at a
fraction of the inference cost. This matters for how the remaining
results in this review should be read: wherever a specific effectiveness
number is cited, it should be understood as a claim made \emph{under
the evaluation conditions of the original study}, not a validated,
generalizable property of the technique. For researchers, this argues
for mandatory leakage audits and cost-normalized comparisons as a
condition of publication in this area; for practitioners, it argues
for re-validating any vendor or paper's effectiveness claims against a
held-out, in-house test set before trusting them in a production
pipeline.

\subsection{Revisiting RQ4: What Remains Open?}
RQ4 is answered most directly by the quality-engineering-safety
synthesis below and by the ten-item research roadmap in
Section~\ref{sec:conclusion}, but two cross-cutting observations are
worth stating plainly here. First, the ``Human factors: trust and
acceptance'' and ``Security of coding and testing agents'' categories
are, jointly, the largest single block of the corpus (30 of 149
studies), which indicates that the research community itself already
recognizes trust and security as the dominant open problem, even though
tooling and standards have not caught up with that recognition.
Second, almost none of the reviewed studies propose or evaluate a
mechanism for \emph{detecting} when an autonomous testing agent has
gone wrong in a way that is not immediately visible as a test failure
(for example, a self-healing agent that has quietly relaxed an
assertion, or a judge model that shares a blind spot with the
generator it is judging). This asymmetry -- extensive work on making
agents capable, comparatively little on making their failures legible
-- is the throughline connecting all four quality-engineering-safety
threads discussed next.

\subsection{Implications for Practitioners and Tool Vendors}
\label{subsec:implications}
For teams evaluating these techniques for production adoption, three
practical implications follow from the synthesis presented in
Section~\ref{sec:results} and Section~\ref{sec:qe-safety}. First, treat
any claim of autonomous operation as scoped to a specific sub-task
(e.g., ``autonomous locator repair'') rather than as a property of an
entire pipeline, and verify the autonomy level against the L0--L5
framing discussed in Section~\ref{sec:conclusion} before relying on it
for unsupervised operation. Second, demand leakage-free, in-house
validation of any effectiveness claim before adoption, given the
documented frequency with which benchmark-reported numbers fail to
reproduce under corrected evaluation. Third, budget for governance
and security from the outset: indirect prompt injection against coding
and testing agents is a demonstrated, quantified threat (Section~\ref{sec:results}),
and human review catches only a fraction of agent-introduced defects
in practice, so relying on manual code review as a sole defense is
untenable.

\subsection{Quality Engineering Safety}
\label{sec:qe-safety}
A cross-cutting theme emerges across four otherwise distinct bodies of
evidence gathered for RQ4: security of testing/coding agents
\cite{zhan2024injecagent,debenedetti2024agentdojo,guo2024redcode,
she2025queryipi,sakib2026trustverify,kozak2025developeraid,
shi2025progent,wu2024isolategpt,wu2026skillshield,xie2026vibeguard,
kokkonis2025rosa}; test-oracle independence and the circularity of AI
testing AI-written code
\cite{bodicoat2026oracle,molina2024oracleautomation,xu2025candor,
wataoka2024selfpreference,panickssery2024evaluators,huang2023selfcorrect,
olausson2023selfrepair}; self-healing/flaky-test stability
\cite{alshammari2024flakyclassifiers,parry2023cannier,
leinen2026detectedundetected,parry2025systemic,lam2020lifecycle,
meta2020pfs,joseph2026zerocost,rohamo2024robotframework,ijam2025healenium};
and governance/accountability for highly autonomous quality engineering
\cite{iso2020tr29119,iso2025ts42119,feng2025autonomylevels,
treude2026accountable,cheng2026threepillar,ibm2024trew,
passi2025overreliance,ibrahim2025overreliance,kiran2026humanai}.

Reading these four threads together reveals a shared structure. In each
case, the core risk is not that an agent fails to perform a task, but
that it performs the task in a way that generates false confidence: an
oracle that quietly accepts subtle bugs because the judge model shares
pretraining blind spots with the generator; a self-healing test that
passes again by modifying a locator to point to the wrong element or by
weakening an assertion; a testing agent hijacked by indirect prompt
injection embedded in a third-party library or test fixture; or a
Level~4--5 autonomous pipeline making release-sign-off decisions without
traceable human accountability. Classical software testing research
relied on an implicit assumption of an independent, trusted, and
adversarially neutral tester; autonomous quality engineering violates
all three assumptions simultaneously.

We argue that resolving this cluster requires treating \emph{quality
engineering safety} as a distinct sub-discipline with its own
mechanisms, rather than as a set of disconnected caveats. Four
candidate mechanisms following this principle are proposed in
Section~\ref{sec:safety-framework} and the supplementary material:
a layered Trust Boundary for Test Agents (TBTA) for the security thread,
Provenance-Partitioned Oracle Verification (PPOV) for the oracle-
circularity thread, Mutation-Gated Dual-Signal Heal Verification
(MG-DSV) for the self-healing thread, and an Autonomy-Tiered Quality
Sign-off Protocol (ATQSP) for the governance thread. None of these four
mechanisms has yet been implemented or empirically evaluated in this
review; they are synthesized proposals for future work, not validated
contributions, and should be read as scoping the research agenda in
Section~\ref{sec:conclusion} rather than as established results.

A further observation, not explored in depth by any single reviewed
study but implied by reading all four threads together, is that the
``second independent signal'' principle has a recursive cost: verifying
a verifier is itself a quality-engineering problem. TBTA's sandboxing
layer must itself be monitored for escape; PPOV's provenance-partitioned
oracle must itself avoid inheriting the generator's blind spots through
shared pretraining data; MG-DSV's mutation-adequacy gate must itself be
protected from the same metric-gaming pressure that motivated the
``Benchmarks and data contamination'' findings under RQ3; and ATQSP's
sign-off protocol is only as trustworthy as the audit log it depends
on. None of the four mechanisms, as currently formalized, fully closes this
regress, which is itself evidence for treating ``verification of
verifiers'' as a standing research problem rather than a one-time
design fix -- a point made explicit in open problem~9 of the roadmap in
Section~\ref{sec:conclusion}.

\subsection{Toward a Testable Safety Framework: Safety-Assured Agentic Quality Engineering}
\label{sec:safety-framework}
The ATQSP sketch introduced in Section~\ref{sec:qe-safety} answers \emph{what} a governance artifact for
autonomous test agents should record, but it does not yet specify
\emph{how} an organization should decide, in a principled and
revisable way, how much autonomy a given agent has earned in a given
context. We elaborate ATQSP into a fuller framework, Safety-Assured
Agentic Quality Engineering (SAAQE), that borrows its central
metaphor from drug-safety practice: an agent, like a drug, is granted
a privilege (dosage, autonomy level) only as far as its accumulated,
monitored track record justifies, and that privilege is continuously
re-evaluated rather than granted once and forgotten. Framed this way,
the four mechanisms proposed in Section~\ref{sec:qe-safety} (TBTA, PPOV,
MG-DSV, and ATQSP) become the \emph{safety guardrails} that feed evidence into a
single accountability substrate, summarized in
Table~\ref{tab:saaqe-components}.

\begin{table}[t]
\caption{Components of the proposed Safety-Assured Agentic Quality
Engineering (SAAQE) framework.}
\label{tab:saaqe-components}
\begin{tabular}{p{0.3\linewidth}p{0.6\linewidth}}
\toprule
Component & What it does \\
\midrule
Safety guardrails &
Five checks: oracle independence from the code under test (PPOV),
mutation-gated self-healing safety (MG-DSV), a sandboxed security
boundary (TBTA), risk-based mandatory human review, and explicit
accountability assignment for release decisions (ATQSP). \\
Knowledge graph &
Traces each escaped defect back through the code and tests to the
specific agent(s) that wrote, healed, or approved them, using
commit-level provenance and bug-inducing-change identification. \\
Reliability score &
A Bayesian score per agent and per task context: clean work raises
the score, escaped defects lower it in proportion to severity, and a
confidence interval reflects the amount of accumulated evidence. \\
Safety boundary (``therapeutic window'') &
The autonomy range in which effort saved by the agent remains high
while the escaped-defect rate stays within an organization-defined
acceptable bound. \\
Privileging &
An agent is granted autonomy level L0--L5
\cite{feng2025autonomylevels} only once its reliability score clears
that level's threshold, and is automatically demoted a level when the
score falls below it -- operationalizing ATQSP's Autonomy Certificate
as a continuously re-evaluated, rather than one-time, grant. \\
Phased rollout and monitoring &
Sandbox evaluation, then shadow mode (the agent acts but a human
retains the decision), then limited production rollout, then full
production, with continuous monitoring of the escaped-defect rate at
each stage before advancing. \\
\bottomrule
\end{tabular}
\end{table}

Read together with the roadmap items in Section~\ref{sec:conclusion}, SAAQE's intended
contribution is fourfold: (i) the framework itself, unifying the four
guardrail mechanisms under one accountability substrate; (ii) a small
set of safety metrics -- escaped-defect rate, masked-regression rate
(valid test behavior silently weakened by a heal or an oracle change),
and attack-success rate under adversarial prompt injection -- that
make ``safety'' measurable rather than aspirational; (iii) a
benchmark suite assembled from existing public datasets (proposed
next); and (iv) an open-source reference implementation of the
knowledge graph and reliability-scoring components, so that the
scoring mechanism itself is auditable rather than a proprietary black
box. We position this contribution within the niche of
\emph{safety-assured agentic quality engineering}: unlike the
capability-focused benchmarks surveyed under RQ1--RQ3, SAAQE is
explicitly a safety-governance layer that sits above whichever
test-generation or self-healing technique an organization already
uses, and is therefore complementary to, not competitive with, the
techniques cataloged in the taxonomy (Section~\ref{sec:results}).

\subsection{An Empirical Validation Plan Using Public Datasets}
\label{sec:saaqe-validation}
Like the four guardrail mechanisms it builds on, SAAQE is a proposed,
not yet empirically validated, mechanism; this subsection outlines how
it \emph{could} be validated before any claim of effectiveness is
made, using data that -- unlike an individual organization's private
CI/CD and defect-tracker history -- is public and therefore
independently reproducible. The guiding idea is to \emph{replay
history}: score agents retrospectively on past work using only
information available at the time, and test whether the resulting
score predicts defects that escaped and were discovered later.
Table~\ref{tab:saaqe-data} summarizes candidate public datasets and
the specific SAAQE component each would exercise.

\begin{table}[t]
\caption{Candidate public datasets for validating the proposed SAAQE
framework.}
\label{tab:saaqe-data}
\begin{tabular}{p{0.26\linewidth}p{0.3\linewidth}p{0.34\linewidth}}
\toprule
Data & What it gives & Use it to test \\
\midrule
AIDev \cite{li2026aidev} &
932{,}791 real agent-authored pull requests from five coding agents
across 116{,}211 repositories &
Reliability scores and the knowledge graph on real, in-the-wild agent
work \\
GitHub issues and bug-fix commits traced with SZZ
\cite{sliwerski2005szz} &
Which earlier commits introduced each later-reported bug &
Escaped defects (bugs found after merge), linked back to the
responsible agent \\
SWE-bench \cite{jimenez2024swebench}, SWT-Bench \cite{slr0282024},
TestGenEval \cite{jain2024testgeneval} &
Agent results on real GitHub issues and tests &
Score calibration; oracle-independence guardrail (PPOV) \\
Defects4J \cite{just2014defects4j}, BugsInPy
\cite{widyasari2020bugsinpy} &
Known, real, previously-fixed faults (prefer the newer, less
contamination-prone releases) &
Whether agent-written tests catch known faults; the masked-regression
rate \\
Flaky-test datasets (e.g., as used by FlakyLens
\cite{slr0442025} and FlakyGuard \cite{slr0492025}) &
Real flaky tests and their fixes &
Heal-safety guardrail (MG-DSV) \\
AgentDojo \cite{debenedetti2024agentdojo}, RepoGuardBench
\cite{li2026repoguardbench} &
Prompt-injection attacks against coding/testing agents &
Security-sandbox guardrail (TBTA) \\
Organization-internal CI/CD and defect-tracker history (where access
permits) &
Production escaped defects with real severity labels &
The strongest validation, as an industrial case study \\
\bottomrule
\end{tabular}
\end{table}

A four-step validation protocol follows from this data: (1)
\emph{build the knowledge graph} from AIDev by linking each agent's
pull requests to the files, tests, and later bug-fix commits
identified via SZZ; (2) \emph{replay the reliability score in time
order}, updating each agent's score at each historical time step using
only evidence that would have been available at that time, to avoid
leaking future information into a retrospective score; (3)
\emph{check predictive validity}, measuring with AUC, precision, and
recall whether agents whose replayed score was low at time $t$
produced more escaped defects in the subsequent period than agents
whose score was high; and (4) \emph{test the safety boundary
counterfactually}, estimating how many of the actually-escaped defects
in the historical record would have been prevented, and how much
measured productivity would have been forgone, had autonomy been
restricted to agents whose score cleared the proposed threshold at
the time. A successful outcome on all four steps would support three
claims: \emph{validity} (the reliability score predicts future
escaped defects), \emph{value} (the safety boundary reduces escaped
defects at an acceptable productivity cost), and \emph{generalization}
(the relationship holds across repositories, languages, and agents).
Three limitations should be stated up front, consistent with this
review's broader emphasis on methodological caution (RQ3,
Section~\ref{sec:threats}): agent-authored pull requests are a recent
phenomenon, so the available history for replay is short;
SZZ-identified bug-inducing commits are known to include labeling
errors, so any SZZ-derived escaped-defect count should be reported
alongside an estimate of SZZ's own error rate rather than treated as
ground truth; and public repositories rarely carry official severity
ratings, so severity must be estimated from proxies such as issue
labels and reporter language, a construct-validity limitation that
should be disclosed rather than hidden behind a single aggregate
metric.
